%% ocean_fcl_routing.tex
%% Ocean FCL Multi-Commodity Routing — MVP Formal Model
%% Phase 1 — must be approved before any code starts
%% Status: Draft v2 — 2026-05-10

\documentclass[11pt,a4paper]{article}

\usepackage{amsmath,amssymb,amsthm}
\usepackage{booktabs}
\usepackage{geometry}
\usepackage{hyperref}
\usepackage{array}
\usepackage{xcolor}
\usepackage{enumitem}
\usepackage{longtable}
\usepackage{graphicx}

\geometry{margin=2.5cm}
\setlength{\parskip}{0.4em}
\hypersetup{colorlinks=true, linkcolor=blue}

\newcommand{\ceil}[1]{\left\lceil #1 \right\rceil}
\newcommand{\floor}[1]{\left\lfloor #1 \right\rfloor}
\newcommand{\defeq}{\triangleq}

\theoremstyle{definition}
\newtheorem{defn}{Definition}
\newtheorem{remark}{Remark}

\title{\textbf{Ocean FCL Multi-Commodity Routing}\\[0.4em]
\large MVP Formal Model --- Phase 1, v2}
\author{AI Freight Agent}
\date{Draft --- 2026-05-10}

\begin{document}
\maketitle

\begin{abstract}
Formal specification of the optimization model for the ocean FCL multi-commodity routing
problem. The model is a Binary Multi-Commodity Flow problem on a commodity-specific
subgraph derived from the physical freight network
(see \texttt{docs/ocean\_fcl\_multi\_shipment\_graph.drawio}).
Scope: FCL only, FEU+TEU containers, string-based carrier allocation capacity, deterministic
time windows. Probabilistic objectives, hazmat/OOG cargo differentiation, and time-phased
capacity management are explicitly deferred to P1.
\end{abstract}

\tableofcontents
\newpage

%%--------------------------------------------------------------------
\section{Problem Statement}
%%--------------------------------------------------------------------

Given a set of shipment requests (commodities), each defined by an origin location,
destination location, cargo volume, cargo weight, and time constraints, find a
minimum-cost assignment of commodities to paths through the multimodal freight network such that:
\begin{itemize}[noitemsep]
  \item every commodity is delivered on or before its deadline,
  \item no ocean sailing exceeds its vessel slot capacity,
  \item no carrier service string exceeds its contracted forwarder allocation in any period, and
  \item optional per-commodity budget caps are respected.
\end{itemize}

\paragraph{MVP scope.}
\begin{itemize}[noitemsep]
  \item FCL only; LCL deferred (PRD Open Question 8).
  \item Container types: FEU (40$'$HC, $\approx 76$\,CBM usable, $\approx 26{,}500$\,kg payload) and TEU
        (20$'$, $\approx 33$\,CBM usable, $\approx 24{,}000$\,kg payload). Optimizer selects mix.
  \item Port dwell: fixed port-specific mean. Commodity-specific inspection model deferred to P1.
  \item Carrier allocation: string-level monthly contracted block, with current utilization tracked.
        Time-phased capacity release deferred to P1.
  \item Objective: minimize total freight cost subject to deterministic time-window feasibility.
        Probabilistic objectives deferred to P1.
  \item No transshipment in MVP. Single-leg ocean sailing only.
\end{itemize}

%%--------------------------------------------------------------------
\section{Graph Structure}
%%--------------------------------------------------------------------

\subsection{Physical Network Graph \texorpdfstring{$G(N, A)$}{G(N,A)}}

The freight network is a directed graph $G = (N, A)$.

\[
  N = N_O \cup N_{\text{POL}} \cup N_{\text{POD},A} \cup N_{\text{POD},E} \cup N_D
\]

\begin{description}[leftmargin=3.2cm, labelwidth=3cm, style=nextline]
  \item[$N_O$ — Origins] One per shipper origin door. Commodity source nodes.
  \item[$N_{\text{POL}}$ — Ports of Loading] Container terminals (e.g., SHA, NGB, SZX). Shipments
        arrive by pre-carriage, depart on vessel sailings.
  \item[$N_{\text{POD},A}$ — POD Arrival] One per discharge port (e.g., LAX$_A$, LGB$_A$, SEA$_A$).
        Vessel arrives; container cannot depart until dwell arc traversed.
  \item[$N_{\text{POD},E}$ — POD Exit] One per discharge port (LAX$_E$, LGB$_E$, SEA$_E$).
        Customs cleared; container available for inland pickup.
  \item[$N_D$ — Destinations] One per delivery door. Commodity sink nodes.
\end{description}

\[
  A = A_{\text{PC}} \cup A_{\text{OC}} \cup A_{\text{DW}} \cup A_{\text{IL}}
\]

\begin{description}[leftmargin=3.2cm, labelwidth=3cm, style=nextline]
  \item[$A_{\text{PC}}$ — Pre-carriage] $(i,j)$, $i \in N_O$, $j \in N_{\text{POL}}$.
        Truck from origin door to container terminal.
  \item[$A_{\text{OC}}$ — Ocean sailings] $(i,j)$, $i \in N_{\text{POL}}$, $j \in N_{\text{POD},A}$.
        One arc per scheduled carrier sailing (specific vessel, departure date, trade lane).
  \item[$A_{\text{DW}}$ — Dwell] $(p_A, p_E)$, one per POD. Vessel unloading plus customs clearance.
  \item[$A_{\text{IL}}$ — Inland] $(p_E, h)$, $p_E \in N_{\text{POD},E}$, $h \in N_D$.
        Truck or intermodal rail, POD exit to destination door.
\end{description}

\begin{remark}[POD node expansion]
Each physical POD is split into two nodes — arrival ($p_A$) and exit ($p_E$) — connected by a
dwell arc. This makes port processing time an explicit model element and provides the structural
hook for the P1 commodity-specific customs inspection model (Section~\ref{sec:deferred}).
\end{remark}

  \begin{figure}[h]
    \centering
    \includegraphics[width=\textwidth]{../docs/shipment_graph.pdf}
    \caption{Multi-shipment freight network $G(N, A)$.}
    \label{fig:multi_shipment_graph}
  \end{figure}

\subsection{Commodity-Specific Subgraph \texorpdfstring{$G(N_k, A_k)$}{G(Nk,Ak)}}
\label{sec:subgraph}

The MILP for commodity $k$ operates on a subgraph $G(N_k, A_k) \subseteq G(N, A)$ containing
only arcs that (i) are time-feasible for commodity $k$ and (ii) lie on at least one complete
path from $o(k)$ to $d(k)$. See Section~\ref{sec:prefilter} for the construction algorithm.




%%--------------------------------------------------------------------
\section{Sets and Indices}
%%--------------------------------------------------------------------

\begin{center}
\begin{tabular}{lp{11cm}}
\toprule
Symbol & Description \\
\midrule
$K$ & Set of commodities, indexed by $k$ \\
$N,\, A$ & All nodes and arcs \\
$A_{\text{PC}},\, A_{\text{OC}},\, A_{\text{DW}},\, A_{\text{IL}}$ & Arc subsets by type \\
$A^k \subseteq A$ & Feasible arc set for commodity $k$ (Section~\ref{sec:prefilter}) \\
$A^k_S \defeq A^k \cap A_S$ & Feasible arcs of type $S$ for commodity $k$ \\
$S$ & Set of carrier service strings (e.g., MSC Tiger, CMA CGM AEX-1) \\
$T$ & Set of allocation time periods (monthly) \\
$P(s)$ & Feasible (POL, POD) port pairs covered by string $s$ \\
$N_k$ & Nodes incident to $A^k$: $N_k \defeq \{n : \exists\,(n,j) \in A^k \text{ or } \exists\,(i,n) \in A^k\}$. Used to eliminate vacuous flow constraints in P.1. \\
$\tau_k(i)$ & Earliest arrival time of commodity $k$ at POL $i$: $\tau_k(i) \defeq t_k^{\text{rdy}} + \mu_{o(k),i}^{\text{PC}}$. Defined here for reference; computed in subgraph pre-filter step 1 (Section~\ref{sec:prefilter}). \\
\bottomrule
\end{tabular}
\end{center}

%%--------------------------------------------------------------------
\section{Parameters}
%%--------------------------------------------------------------------

\subsection{Commodity Parameters}

\begin{center}
\begin{tabular}{llp{6cm}l}
\toprule
Parameter & Symbol & Description & Unit \\
\midrule
Volume & $v_k$ & Cargo volume & CBM \\
Weight & $w_k$ & Gross weight & kg \\
FEU-equivalent count & $n_k$ & See Eq.~\eqref{eq:nk}; pre-computed & — \\
Cargo ready & $t_k^{\text{rdy}}$ & Earliest pickup from origin door & datetime \\
Deadline & $T_k^{\text{dead}}$ & Latest acceptable arrival at destination & datetime \\
Origin & $o(k)$ & Source node $\in N_O$ & — \\
Destination & $d(k)$ & Sink node $\in N_D$ & — \\
Cargo type & $\tau_k$ & General / hazmat / temp-controlled / OOG & categorical \\
HS code & $\text{hs}_k$ & 6-digit Harmonized System code (P1 customs model) & — \\
Incoterm & $\text{inco}_k$ & EXW / FOB / CIF / DDP & — \\
Budget cap & $B_k$ & Hard cost ceiling; $\infty$ if not set & USD \\
\bottomrule
\end{tabular}
\end{center}

\paragraph{Container count.}
\begin{equation}
  n_k \defeq \max\!\left(\ceil{\frac{v_k}{76}},\ \ceil{\frac{w_k}{26{,}500}}\right)
  \label{eq:nk}
\end{equation}

This accounts for both volume (76\,CBM per 40$'$HC FEU) and weight (26,500\,kg payload limit
per FEU). $n_k$ is the minimum number of FEU-sized moves for land transport; the optimizer
may substitute some FEUs with TEUs on ocean arcs to reduce cost (Section~\ref{sec:containermix}).

\paragraph{Container specifications.}

\begin{center}
\begin{tabular}{llllll}
\toprule
Container type & Code & TEU slots & Usable volume & Payload & Ocean freight rate \\
\midrule
20$'$ standard & TEU & 1 & 33\,CBM & 24{,}000\,kg & $\rho \cdot c_{ij}^{\text{FEU}}$, $\rho$ lane-specific (see \S\ref{sec:ocean-params}) \\
40$'$ standard & FEU (std) & 2 & 67\,CBM & 26{,}500\,kg & $c_{ij}^{\text{FEU}} - \$50$--\$200 vs.\ HC \\
40$'$ High Cube & FEU (HC) & 2 & 76\,CBM & 26{,}500\,kg & $c_{ij}^{\text{FEU}}$ (baseline; defined in \S\ref{sec:ocean-params}) \\
\bottomrule
\end{tabular}
\end{center}

\noindent MVP uses the 40$'$HC as the FEU type (76\,CBM, 26{,}500\,kg). The 40$'$ standard
is deferred to P1. TEU-only routing for heavy, low-volume cargo is also P1.

\paragraph{Container rate ratios.}
The TEU/FEU cost ratio $\rho = c_{ij}^{\text{TEU}} / c_{ij}^{\text{FEU}}$ varies by trade lane.
Historical averages 2020--May 2024 (Xeneta\footnote{\url{https://www.xeneta.com/blog/beware-a-potential-flaw-in-index-linked-ocean-container-shipping-contracts-which-could-cost-millions-of-dollars}}):

\begin{center}
\begin{tabular}{llcc}
\toprule
Trade lane & Route & $\rho$ (avg) & $\rho$ (range) \\
\midrule
TPEB & China -- US West Coast & 0.79 & 0.71--0.86 \\
FEWB & Shanghai -- Antwerp    & 0.56 & 0.52--0.61 \\
TAWB & N.\ Europe -- US East Coast & 0.78 & 0.68--0.94 \\
\bottomrule
\end{tabular}
\end{center}
\noindent Spot ratios fluctuate with equipment balance and demand cycles; on Far East -- North
Europe, $\rho$ fell to 0.58 at peak demand (July 2024).\footnote{\url{https://www.xeneta.com/blog/dramatic-increases-in-feu-and-teu-ocean-container-market-spreads-key-considerations-for-shippers}}
FEU (HC) vs.\ FEU (std) ocean freight premium: \$50--\$200 per container, reflecting the 13\%
volume advantage of the HC over the standard 40$'$.\footnote{\url{https://www.icontainers.com/help/40-foot-high-cube-container/}}
FEU (std) is not separately modelled in MVP.


\begin{remark}[Cargo type in MVP]
$\tau_k$ is recorded for all commodities but is not used in MVP constraints. Hazmat, OOG, and
reefer cargo are routed identically to general cargo in MVP. The instance generator should
restrict generated commodities to general cargo unless explicitly testing cargo-type extensions.
\end{remark}

\subsection{Pre-Carriage Arc Parameters}

For each $(i,j) \in A_{\text{PC}}$:

\begin{center}
\begin{tabular}{llp{6cm}l}
\toprule
Parameter & Symbol & Description & Unit \\
\midrule
Cost & $c_{ij}^{\text{PC}}$ & Truck rate per container move. Scales by $n_k$ (one truck per container). & USD/move \\
Mean transit & $\mu_{ij}^{\text{PC}}$ & Mean days, origin door to terminal gate & days \\
Std dev & $\sigma_{ij}^{\text{PC}}$ & Std dev of transit time & days \\
Road distance & $\text{dist}_{ij}^{\text{PC}}$ & Road distance (informational; used by generator) & km \\
\bottomrule
\end{tabular}
\end{center}

\subsection{Ocean Sailing Arc Parameters}
\label{sec:ocean-params}

For each $(i,j) \in A_{\text{OC}}$:

\begin{center}
\begin{tabular}{llp{5.5cm}l}
\toprule
Parameter & Symbol & Description & Unit \\
\midrule
ETD & $\text{ETD}_{ij}$ & Estimated time of departure from POL $i$ & datetime \\
CY cutoff & $\text{CYC}_{ij}$ & Latest container arrival at terminal. $\text{CYC}_{ij} \leq \text{ETD}_{ij}$; typically $\text{ETD} - 4\,\text{d}$ & datetime \\
Booking cutoff & $\text{BKG}_{ij}$ & Latest booking confirmation with carrier. \textit{[Note C2: Not enforced in MVP. P1: add as a hard pre-filter --- exclude any sailing where $\text{BKG}_{ij}$ has passed at solve time.]} & datetime \\
FEU rate & $r_{ij}^{\text{FEU}}$ & Base ocean freight rate per 40$'$HC container & USD/FEU \\
TEU rate & $r_{ij}^{\text{TEU}}$ & Base rate per 20$'$ container. Typically $0.56$--$0.86 \times r_{ij}^{\text{FEU}}$ (see rate ratio table) & USD/TEU \\
BAF & $\text{BAF}_{ij}$ & Bunker adjustment per container (same fraction applied to FEU and TEU rates) & USD \\
PSS & $\text{PSS}_{ij}$ & Peak Season Surcharge per container; zero if not applicable. \textit{[Note H5: Not in MVP cost formula. P1: add $+\,\text{PSS}_{ij}$ to the $c_{ij}^{\text{FEU}}$ derivation below and to the instance generator.]} & USD \\
THC (POL) & $\text{THC}_{ij}^{\text{POL}}$ & Terminal handling at origin per container & USD \\
THC (POD) & $\text{THC}_{ij}^{\text{POD}}$ & Terminal handling at destination per container & USD \\
Vessel capacity & $\text{cap}_{ij}^{\text{TEU}}$ & Per-sailing slot cap (constraint P.2). MVP proxy: $\alpha \cdot \text{alloc}(s_{ij}, t_{ij})$, $\alpha \approx 0.20$ (configurable). & TEU \\
String & $s_{ij} \in S$ & Carrier service string this sailing belongs to & — \\
Period & $t_{ij} \in T$ & Allocation period (month) this ETD falls in & — \\
Mean transit & $\mu_{ij}^{\text{OC}}$ & Mean transit days, POL to POD & days \\
Std dev & $\sigma_{ij}^{\text{OC}}$ & Std dev of transit & days \\
Carrier & $\text{carrier}_{ij}$ & Carrier name (e.g., MSC, CMA CGM, COSCO) & — \\
Vessel & $\text{vessel}_{ij}$ & Vessel name & — \\
\bottomrule
\end{tabular}
\end{center}

\noindent All-in per-FEU cost (derived):
\[
  c_{ij}^{\text{FEU}} \defeq r_{ij}^{\text{FEU}} + \text{BAF}_{ij} + \text{THC}_{ij}^{\text{POL}} + \text{THC}_{ij}^{\text{POD}}
\]
Similarly $c_{ij}^{\text{TEU}}$ uses $r_{ij}^{\text{TEU}}$ in place of $r_{ij}^{\text{FEU}}$.

\subsection{Dwell Arc Parameters}

For each $(p_A, p_E) \in A_{\text{DW}}$:

\begin{center}
\begin{tabular}{llp{5cm}l}
\toprule
Parameter & Symbol & Description & MVP value \\
\midrule
Mean dwell & $\mu_p^{\text{DW}}$ & Unloading + customs clearance (no exam) & USLAX/USLGB: 3.5\,d; USSEA: 2.5\,d \\
Std dev & $\sigma_p^{\text{DW}}$ & Std dev of total dwell & USLAX/USLGB: 1.5\,d; USSEA: 1.0\,d \\
Free time & $f_p$ & Free days before storage charges & 5\,d \\
\bottomrule
\end{tabular}
\end{center}

\noindent P1: dwell becomes a function of $\text{hs}_k$, C-TPAT status, origin country, consignee
history. See Section~\ref{sec:deferred}.

\subsection{Inland Arc Parameters}

For each $(p_E, h) \in A_{\text{IL}}$:

\begin{center}
\begin{tabular}{llp{5.5cm}l}
\toprule
Parameter & Symbol & Description & Unit \\
\midrule
Cost & $c_{p_E h}^{\text{IL}}$ & Flat rate per container move (one chassis per container) & USD/move \\
Mean transit & $\mu_{p_E h}^{\text{IL}}$ & Mean days, POD exit to destination & days \\
Std dev & $\sigma_{p_E h}^{\text{IL}}$ & Std dev of transit & days \\
Mode & $\text{mode}_{p_E h}$ & FTL truck or intermodal rail & — \\
Road distance & $\text{dist}_{p_E h}^{\text{IL}}$ & Informational & km \\
\bottomrule
\end{tabular}
\end{center}

\subsection{Carrier Allocation Parameters (String-Level)}
\label{sec:alloc-params}

Ocean carriers sell capacity in blocks on \emph{service strings}. A string is a fixed port-call
loop (e.g., MSC Tiger TPEB: SHA $\to$ NGB $\to$ SZX $\to$ USLAX $\to$ USLGB). A forwarder's
contracted allocation covers \emph{all} (POL, POD) pairs on a string — not a specific port pair.

For each string $s \in S$ and period $t \in T$:

\begin{center}
\begin{tabular}{llp{6cm}l}
\toprule
Parameter & Symbol & Description & Unit \\
\midrule
Contracted allocation & $\text{alloc}(s,t)$ & Forwarder's contracted block on string $s$ in period $t$ & TEU \\
Current utilization & $\text{util}(s,t)$ & TEU already booked/confirmed on $s$ in $t$ at solve time & TEU \\
Remaining capacity & $\text{rem}(s,t)$ & $\text{alloc}(s,t) - \text{util}(s,t)$; input to optimizer & TEU \\
Vessel cap fraction & $\alpha$ & Fraction of string allocation used as per-sailing slot cap proxy. Default $0.20$. & — \\
\bottomrule
\end{tabular}
\end{center}

$\text{rem}(s,t)$ is computed before each solve from the shipment state store. The optimizer
sees only remaining capacity — utilization is an external state input, not a decision variable.

\begin{remark}[String vs.\ port-pair capacity]
A forwarder's contracted block on MSC Tiger (e.g., 400\,TEU slots/month) applies to
\emph{all sailings} of that string in the month, across \emph{all} POL/POD combinations.
A booking on SHA$\to$USLAX and a booking on NGB$\to$USLGB both draw from the same pool.
\end{remark}

\begin{remark}[Unit convention]
All allocation quantities --- $\text{alloc}(s,t)$, $\text{util}(s,t)$, $\text{rem}(s,t)$ ---
are in \textbf{TEU slots}, consistent with $\text{slots}(k,ij)$.
Real carrier BSA contracts may be quoted in FEU (common on TPEB) or TEU depending on carrier
and trade lane. Convert to TEU slots at input time: $\text{alloc}_{\text{FEU}} \times 2 =
\text{alloc in TEU slots}$.
\end{remark}

\subsection{Pre-Computed Container Mix}
\label{sec:containermix}

For each commodity $k$ and each feasible ocean arc $(i,j) \in A^k_{\text{OC}}$, we pre-compute
the optimal container mix that minimizes ocean cost while satisfying volume and weight constraints:

\begin{align}
  (f_k^*, t_k^*)\big|_{ij} &= \arg\min_{f,\,t \,\geq\, 0,\; \text{integer}} \;\;
    f \cdot c_{ij}^{\text{FEU}} + t \cdot c_{ij}^{\text{TEU}} \notag \\
  &\text{s.t.}\quad 76f + 33t \geq v_k \quad \text{(volume)} \label{eq:mix-vol} \\
  &\phantom{\text{s.t.}}\quad 26{,}500f + 24{,}000t \geq w_k \quad \text{(weight)} \label{eq:mix-wt} \\
  &\phantom{\text{s.t.}}\quad f, t \in \mathbb{Z}_{\geq 0} \notag
\end{align}

This is a 2-variable integer program, solvable analytically (see Section~\ref{sec:mix-algo}).
The results feed into the main MILP as parameters:

\begin{center}
\begin{tabular}{lll}
\toprule
Derived parameter & Expression & Description \\
\midrule
$\text{cost}(k,ij)$ & $f_k^*\big|_{ij} \cdot c_{ij}^{\text{FEU}} + t_k^*\big|_{ij} \cdot c_{ij}^{\text{TEU}}$ & Optimal ocean cost for commodity $k$ on sailing $(i,j)$ \\
$\text{slots}(k,ij)$ & $2 f_k^*\big|_{ij} + t_k^*\big|_{ij}$ & TEU slots consumed on the vessel \\
\bottomrule
\end{tabular}
\end{center}

\begin{remark}[$n_k$ and physical container count]
$n_k$ counts FEU-sized moves and is used to price the land legs: $n_k$ trucks at rate
$c^{\text{PC}}$ or $c^{\text{IL}}$, one chassis per container. On ocean arcs, the mix
pre-computation may assign a different combination of FEUs and TEUs. The total physical
container count on sailing $(i,j)$ is $f_k^*\big|_{ij} + t_k^*\big|_{ij}$, which is always
$\geq n_k$ (you cannot cover the same cargo with fewer physical containers than the FEU
minimum). The gap is at most 1:
\[
  \bigl(f_k^*\big|_{ij} + t_k^*\big|_{ij}\bigr) - n_k \leq 1
\]
This holds because replacing one FEU with two or more TEUs is never cost-effective at
$\rho \in (0.56, 0.86)$: three TEUs cost $3\rho \geq 1.68\,c_{ij}^{\text{FEU}}$ but cover
only $99\,\text{CBM}$, less volume per dollar than one FEU. The only viable substitution is
replacing the single residual FEU with one TEU when the remaining volume fits ($\leq 33$\,CBM),
which leaves the total count unchanged at $n_k$. In practice the deviation is zero; the bound
$\leq 1$ is a theoretical guarantee.
\end{remark}

\paragraph{Container mix analytical solution.}
\label{sec:mix-algo}
Given cost ratio $\rho = c_{ij}^{\text{TEU}} / c_{ij}^{\text{FEU}} \in (0.56, 0.86)$ (see rate ratio table above) and constraints
\eqref{eq:mix-vol}--\eqref{eq:mix-wt}:

\begin{enumerate}[noitemsep]
  \item Compute the minimum FEU count needed from volume and weight:
        $f_v = \ceil{v_k / 76}$, $f_w = \ceil{w_k / 26{,}500}$. Set $f_{\min} = \max(f_v, f_w)$.
  \item For each candidate $f \in \{f_{\min},\, f_{\min}-1,\, \ldots,\, 0\}$:
        \begin{enumerate}[noitemsep, label=\alph*.]
          \item $t_v = \max\!\left(0,\, \ceil{(v_k - 76f)/33}\right)$
                (minimum TEUs to cover remaining volume).
          \item $t_w = \max\!\left(0,\, \ceil{(w_k - 26{,}500f)/24{,}000}\right)$
                (minimum TEUs to cover remaining weight).
          \item $t = \max(t_v, t_w)$.
          \item Evaluate $C(f,t) = f \cdot c_{ij}^{\text{FEU}} + t \cdot c_{ij}^{\text{TEU}}$.
        \end{enumerate}
  \item Return $(f^*, t^*) = \arg\min C(f,t)$ over all candidates from step 2.
\end{enumerate}

The enumeration runs over at most $f_{\min}+1$ candidates, which is small in practice
(typical $f_{\min} \leq 5$). Since $\rho \in (0.63, 0.83)$, one FEU covers
$76/33 \approx 2.3$ TEU-equivalents of volume at a cost of $1/\rho \in (1.20, 1.59)$
TEU-equivalents --- so FEU is always cheaper per CBM. The optimal solution is therefore
$(f_{\min}, 0)$ or $(f_{\min}-1, t)$ for small residual $t$, but explicit enumeration
guarantees optimality without relying on this structural property.

%%--------------------------------------------------------------------
\section{Commodity-Specific Subgraph Construction}
\label{sec:prefilter}
%%--------------------------------------------------------------------

Before the MILP is built, we construct $G(N_k, A_k)$ for each commodity $k$. The subgraph
contains only arcs that are (i) time-feasible and (ii) reachable on a complete path from
$o(k)$ to $d(k)$. Any arc not on such a path is excluded.

\paragraph{Key structural observation.}
In this network, each commodity $k$ has \emph{at most one pre-carriage arc per POL}: the arc
$(o(k), j)$ for each $j \in N_{\text{POL}}$. The transit time $\mu_{o(k),j}^{\text{PC}}$ is
therefore unambiguous for a given (commodity, POL) pair. This eliminates the arc-pairing
ambiguity that would arise if multiple different pre-carriage routes from different origins
could reach the same POL with different transit times.

\paragraph{Construction algorithm.}

\begin{enumerate}
  \item \textbf{Pre-carriage pass.} For each POL $i \in N_{\text{POL}}$, compute earliest arrival:
    \[
      \tau_k(i) \defeq t_k^{\text{rdy}} + \mu_{o(k),i}^{\text{PC}}
    \]
    Arc $(o(k), i)$ is \emph{candidate} iff $\exists\, (i,j) \in A_{\text{OC}}$ such that
    $\tau_k(i) \leq \text{CYC}_{ij}$ (can make at least one sailing from $i$).

  \item \textbf{Ocean sailing pass.} For each arc $(i,j) \in A_{\text{OC}}$:
    \[
      (i,j) \in A^k_{\text{OC}} \iff
        \tau_k(i) \leq \text{CYC}_{ij}
        \;\text{ and }\;
        \text{ETD}_{ij} + \mu_{ij}^{\text{OC}} + \mu_p^{\text{DW}} + \mu_{p_E,d(k)}^{\text{IL}}
          \leq T_k^{\text{dead}}
    \]
    where $j \in N_{\text{POD},A}$, $p$ is the physical port at $j$, and $p_E$ is its exit node.
    The first condition is the CY cutoff check (directly against commodity $k$'s pre-carriage
    transit to POL $i$ — no ambiguity). The second is the end-to-end deadline check.

  \item \textbf{Dwell pass.} Include dwell arc $(p_A, p_E)$ iff $\exists\, (i,j) \in A^k_{\text{OC}}$
    with $j = p_A$ (i.e., some feasible sailing arrives at that POD).

  \item \textbf{Inland pass.} For each POD exit $p_E$ and inland arc $(p_E, h) \in A_{\text{IL}}$:
    \[
      (p_E, h) \in A^k_{\text{IL}} \iff h = d(k) \;\text{ and }\;
      \exists\, (i,j) \in A^k_{\text{OC}} : \text{ETD}_{ij} + \mu_{ij}^{\text{OC}}
        + \mu_p^{\text{DW}} + \mu_{p_E h}^{\text{IL}} \leq T_k^{\text{dead}}
    \]

  \item \textbf{Trim pre-carriage.} Remove any pre-carriage arc $(o(k), i)$ for which no
    ocean arc $(i,j) \in A^k_{\text{OC}}$ remains adjacent to $i$.

  \item \textbf{Reachability sweep.} Run a forward BFS from $o(k)$ and a backward BFS from
    $d(k)$ on the candidate arc set. Retain only arcs on paths reachable from both ends.
    Result: $G(N_k, A_k)$ where every arc lies on at least one complete $o(k) \to d(k)$ path.
\end{enumerate}

\paragraph{Pre-filter scope guard.}
\textit{[Note H3: The pre-filter only processes (origin, destination) pairs whose trade lane
is covered by the configured carrier services. O-D pairs outside the configured scope must be
hard-rejected before the pre-filter runs, with a structured error returned to the agent.]}

\textit{[Note H4: Hazmat, OOG, and reefer commodities ($\tau_k \neq \text{general}$) must be
rejected at the pre-filter boundary with a structured infeasibility report. Routing them
identically to general cargo produces silently incorrect results. The instance generator
already restricts to general cargo (see §4.1 Remark); enforce the same check at runtime.]}

\textit{[Note H2: The deadline check in step 2 uses mean transit times. When a sailing's margin
is tight, mean-based inclusion overstates on-time probability. P1: add a safety buffer
$\delta \approx 1$--$2\,\text{d}$ to the deadline check:
$\text{ETD}_{ij} + \mu_{ij}^{\text{OC}} + \mu_p^{\text{DW}} + \mu_{p_E,d(k)}^{\text{IL}} + \delta \leq T_k^{\text{dead}}$.
The CYC risk note in Open Item 3 partially addresses this for the pre-carriage side.]}

\paragraph{Infeasible commodities.}
If $A^k = \emptyset$ after the above (no feasible path exists), commodity $k$ is \emph{removed
from $K$ before the MILP is built}. The pre-solve returns a structured infeasibility report:
\[
  \texttt{\{commodity\_id, reason: "no\_feasible\_path",
  earliest\_achievable\_days, deadline\_shortfall\_days\}}
\]
The agent surfaces this to the user with a recommended recovery action (relax deadline,
consider air mode). The MILP is never passed an infeasible commodity.

%%--------------------------------------------------------------------
\section{Decision Variables}
%%--------------------------------------------------------------------

\[
  x_{ij}^k \in \{0, 1\} \qquad \forall\, k \in K,\; (i,j) \in A^k
\]

$x_{ij}^k = 1$ if commodity $k$ is routed on arc $(i,j)$.

\begin{remark}[No-split semantics]
$x_{ij}^k = 1$ means \emph{all} containers for commodity $k$ traverse arc $(i,j)$ together.
Commodity splitting across multiple sailings or paths is prohibited. For an ocean arc, this
means all $f_k^*\big|_{ij}$ FEUs and $t_k^*\big|_{ij}$ TEUs depart on the same vessel.
This reflects standard FCL booking practice: one booking covers all containers for a shipment.
\end{remark}

%%--------------------------------------------------------------------
\section{Objective Function}
%%--------------------------------------------------------------------

\begin{equation}
  \min \;\; Z = \sum_{k \in K} \left[
    \underbrace{\sum_{(i,j) \in A^k_{\text{PC}}} n_k \cdot c_{ij}^{\text{PC}} \cdot x_{ij}^k}_{\text{pre-carriage}}
    +
    \underbrace{\sum_{(i,j) \in A^k_{\text{OC}}} \text{cost}(k,ij) \cdot x_{ij}^k}_{\text{ocean (optimal mix)}}
    +
    \underbrace{\sum_{(p_E,h) \in A^k_{\text{IL}}} n_k \cdot c_{p_E h}^{\text{IL}} \cdot x_{p_E h}^k}_{\text{inland}}
  \right]
  \label{eq:obj}
\end{equation}

\paragraph{Cost structure.}
\begin{itemize}[noitemsep]
  \item \textit{Pre-carriage:} $n_k$ trucks (one chassis per container), each at flat rate $c_{ij}^{\text{PC}}$.
  \item \textit{Ocean:} $\text{cost}(k,ij)$ is the pre-computed optimal cost for commodity $k$ on sailing $(i,j)$,
        accounting for the optimal FEU/TEU mix. Dwell cost (port storage beyond free time) deferred to P1.
  \item \textit{Inland:} $n_k$ chassis moves, each at flat rate $c_{p_E h}^{\text{IL}}$ (one chassis per container for FCL drayage).
\end{itemize}

\textit{[Note C1: Pre-carriage and inland costs use $n_k$ as the container-count multiplier. The actual count on the chosen sailing is $f_k^*\big|_{ij} + t_k^*\big|_{ij}$. Because $2\rho > 1$ for all $\rho \in (0.56, 0.86)$, the optimizer never replaces a FEU with 2 TEUs, so $f^*+t^* = n_k$ in every feasible solution. Using $n_k$ is therefore exact, not an approximation.]}

%%--------------------------------------------------------------------
\section{Constraints}
%%--------------------------------------------------------------------

\subsection{Flow Conservation}

For each commodity $k \in K$ and node $n \in N_k$:

\begin{equation}
  \sum_{(n,j) \in A^k} x_{nj}^k \;-\; \sum_{(i,n) \in A^k} x_{in}^k \;=\; b_k(n)
  \qquad \forall\, k \in K,\; n \in N_k
  \label{eq:flow}
\end{equation}

\[
  b_k(n) = \begin{cases}
    +1 & n = o(k) \quad \text{(net outflow = 1 at source)} \\
    -1 & n = d(k) \quad \text{(net inflow = 1 at sink)} \\
    \phantom{+}0 & \text{otherwise}
  \end{cases}
\]

Flow conservation ensures each commodity routes along exactly one complete path (one arc of
each type: PC, OC, DW, IL). The subgraph construction guarantees any arc in $A^k$ is on a
valid complete path, so feasibility is structurally enforced — no additional time constraints
are needed in the MILP.

\subsection{Vessel Capacity Constraint}

For each ocean arc $(i,j) \in A_{\text{OC}}$:

\begin{equation}
  \sum_{k \in K:\, (i,j) \in A^k_{\text{OC}}} \text{slots}(k,ij) \cdot x_{ij}^k
  \;\leq\; \text{cap}_{ij}^{\text{TEU}}
  \qquad \forall\, (i,j) \in A_{\text{OC}}
  \label{eq:vcap}
\end{equation}

The total TEU slots consumed on a single sailing cannot exceed that sailing's slot
availability. Without this constraint, the string allocation constraint alone allows all
monthly forwarder allocation to concentrate on one departure, which is physically impossible.

\paragraph{Proxy for $\text{cap}_{ij}^{\text{TEU}}$.}
Real per-sailing slot availability is not directly observable by a forwarder. As a
non-binding placeholder:
\[
  \text{cap}_{ij}^{\text{TEU}} \defeq \alpha \cdot \text{alloc}(s_{ij},\, t_{ij})
\]
$\alpha \approx 0.20$ (configurable). Rationale: with $\approx 4$ sailings per month per
string, $\alpha = 0.20$ caps any single departure at roughly 80\,\% of the full monthly
contracted block --- large enough to be non-binding on typical MVP instances, tight enough
to prevent degenerate concentration. See Open Item~7 for the path to a real estimate.

\subsection{String Allocation Constraint}

For each carrier service string $s \in S$ and time period $t \in T$:

\begin{equation}
  \sum_{\substack{(i,j) \in A_{\text{OC}} \\ s_{ij} = s,\; t_{ij} = t}}
  \sum_{k \in K:\, (i,j) \in A^k_{\text{OC}}} \text{slots}(k,ij) \cdot x_{ij}^k
  \;\leq\; \text{rem}(s,t)
  \qquad \forall\, s \in S,\; t \in T
  \label{eq:alloc}
\end{equation}

The total TEU slots consumed across \emph{all sailings of string $s$ in period $t$} cannot
exceed the forwarder's remaining allocated capacity. This constraint couples routing decisions
across all (POL, POD) pairs on the same string within the same month.


\subsection{Budget Cap (Optional)}

For each $k \in K$ with finite $B_k$:

\begin{equation}
  \sum_{(i,j) \in A^k_{\text{PC}}} n_k c_{ij}^{\text{PC}} x_{ij}^k
  + \sum_{(i,j) \in A^k_{\text{OC}}} \text{cost}(k,ij)\, x_{ij}^k
  + \sum_{(p_E,h) \in A^k_{\text{IL}}} n_k c_{p_E h}^{\text{IL}} x_{p_E h}^k
  \leq B_k
  \label{eq:budget}
\end{equation}

\subsection{Variable Domains}

\begin{equation}
  x_{ij}^k \in \{0,1\} \qquad \forall\, k \in K,\; (i,j) \in A^k \label{eq:domain}
\end{equation}

%%--------------------------------------------------------------------
\section{Complete Formulation}
%%--------------------------------------------------------------------

\begin{align}
  \min \quad & Z = \sum_{k \in K} \left[
    \sum_{(i,j) \in A^k_{\text{PC}}} n_k c_{ij}^{\text{PC}} x_{ij}^k
    + \sum_{(i,j) \in A^k_{\text{OC}}} \text{cost}(k,ij)\, x_{ij}^k
    + \sum_{(p_E,h) \in A^k_{\text{IL}}} n_k c_{p_E h}^{\text{IL}} x_{p_E h}^k
  \right] \tag{P} \label{eq:P} \\[4pt]
  \text{s.t.} \quad
  & \sum_{(n,j) \in A^k} x_{nj}^k - \sum_{(i,n) \in A^k} x_{in}^k = b_k(n)
    && \forall\, k \in K,\; n \in N_k \tag{P.1} \\[2pt]
  & \sum_{k \in K:\, (i,j) \in A^k_{\text{OC}}} \text{slots}(k,ij)\, x_{ij}^k
      \leq \text{cap}_{ij}^{\text{TEU}}
    && \forall\, (i,j) \in A_{\text{OC}} \tag{P.2} \\[2pt]
  & \sum_{\substack{(i,j) \in A_{\text{OC}} \\ s_{ij}=s,\, t_{ij}=t}} \sum_{k:\, (i,j) \in A^k_{\text{OC}}} \text{slots}(k,ij)\, x_{ij}^k
      \leq \text{rem}(s,t)
    && \forall\, s \in S,\; t \in T \tag{P.3} \\[2pt]
  & \text{(P.4): budget constraints \eqref{eq:budget} for commodities with finite } B_k \notag \\[2pt]
  & x_{ij}^k \in \{0,1\}
    && \forall\, k \in K,\; (i,j) \in A^k \tag{P.5}
\end{align}

\begin{remark}[Problem structure]
(P) is a Binary Multi-Commodity Network Flow (BMCNF). The LP relaxation of a pure MCF is
TU; coupling via the string allocation constraint (P.3) breaks TU, making
integrality non-guaranteed. For small MVP instances ($|K| \leq 50$), HiGHS solves to
optimality easily. At forwarder scale ($|K| \geq 500$), LP bound quality and symmetry
breaking become relevant (see Section~\ref{sec:solver}).
\end{remark}

%%--------------------------------------------------------------------
\section{Graph Decomposition for Independent Subproblems}
\label{sec:decomp}
%%--------------------------------------------------------------------

When routing a mixed batch of commodities (e.g., from both TPEB and FEWB trade lanes),
groups of commodities may share no supply arcs and no allocation pools — making their
routing decisions completely independent. Identifying and solving these independently
reduces MILP solve time and enables parallelism.

\subsection{Commodity-Supply Graph}

Define a commodity-supply graph $H = (K, E_H)$:
\begin{itemize}[noitemsep]
  \item \textbf{Nodes:} one per commodity $k \in K$.
  \item \textbf{Edges:} $(k_1, k_2) \in E_H$ iff commodities $k_1$ and $k_2$ share supply:
    \begin{enumerate}[noitemsep]
      \item $\exists\, (i,j) \in A^{k_1}_{\text{OC}} \cap A^{k_2}_{\text{OC}}$ (same feasible sailing), \textit{or}
      \item $\exists\, s \in S,\; t \in T$ such that both $k_1$ and $k_2$ have a feasible
            sailing on string $s$ in period $t$: i.e., $\exists\,(i,j) \in A^{k_1}_{\text{OC}}$
            with $s_{ij}=s,\; t_{ij}=t$, and $\exists\,(i',j') \in A^{k_2}_{\text{OC}}$
            with $s_{i'j'}=s,\; t_{i'j'}=t$ (same allocation pool). Arc membership in
            $A^k_{\text{OC}}$ is determined by the subgraph pre-filter
            (Section~\ref{sec:prefilter}), not the MILP solve itself.
    \end{enumerate}
\end{itemize}

\subsection{Decomposition Algorithm}

\begin{enumerate}[noitemsep]
  \item Build $H$ from $\{A^k_{\text{OC}}\}_{k \in K}$ and the sailing-to-string mapping $s_{ij}$.
  \item Find connected components $\mathcal{C}_1, \mathcal{C}_2, \ldots$ of $H$.
  \item For each component $\mathcal{C}_i$: solve the MILP restricted to $K_i = \mathcal{C}_i$.
  \item Merge solutions.
\end{enumerate}

\textit{[Note M9: $\text{rem}(s,t)$ must be computed once before the decomposition and
held fixed for all component solves. Updating it between components --- e.g., subtracting
slots used by a solved component before solving the next --- is incorrect: the allocation
constraint in each component references the full remaining capacity at solve time, and
components are, by construction, independent.]}

\paragraph{Example.} TPEB commodities (SHA$\to$USLAX on MSC Tiger) and FEWB commodities
(SHA$\to$NLRTM on MSC Silk) use different strings with separate allocation pools and no
overlapping sailings. They form disconnected components in $H$ and are solved independently.
Within TPEB, two commodities that both have SHA$\to$USLAX in their feasible arc sets are
connected in $H$ (shared sailing, shared allocation) and solved jointly.

%%--------------------------------------------------------------------
\section{Transit Time Estimation}
\label{sec:transit}
%%--------------------------------------------------------------------

Transit times are used in the subgraph pre-filter and by the instance generator. They are
parameters in MVP; in P1 they become distribution parameters for chance constraints.

\subsection{Trucking (Pre-Carriage and Inland)}

\paragraph{Geodesic distance (Haversine):}
\[
  d_{ij}^{\text{geo}} = 2R \arcsin\!\sqrt{
    \sin^2\tfrac{\phi_j - \phi_i}{2}
    + \cos\phi_i \cos\phi_j \sin^2\tfrac{\lambda_j - \lambda_i}{2}
  }, \quad R = 6{,}371\,\text{km}
\]

\paragraph{Road distance:}
\[
  \text{dist}_{ij}^{\text{road}} = \kappa \cdot d_{ij}^{\text{geo}}, \quad
  \kappa = \begin{cases}
    1.25 & \text{China pre-carriage} \\
    1.20 & \text{US inland}
  \end{cases}
\]

\paragraph{Transit time:}
\[
  \mu_{ij}^{\text{truck}} = \frac{\text{dist}_{ij}^{\text{road}}}{v^{\text{truck}}}, \quad
  v^{\text{truck}} = \begin{cases}
    600\,\text{km/day} & \text{China (urban congestion, regulations)} \\
    800\,\text{km/day} & \text{US OTR (HOS 11-hour limit)}
  \end{cases}
\]

\paragraph{Std dev:} $\sigma_{ij}^{\text{truck}} = 0.15\,\mu_{ij}^{\text{truck}}$ (15\% CV).

\paragraph{Transit time range (for robust model, P1):}
Each arc carries $[\mu - 2\sigma,\; \mu + 2\sigma]$ as a 95\% interval. The instance
generator outputs both mean and std dev; the range is the hook for P1 robustness.

\subsection{Ocean Sailing}

\paragraph{Sailing distance:}
\[
  \text{dist}_{ij}^{\text{sail}} = 1.15 \cdot d_{ij}^{\text{geo}} \quad \text{(Trans-Pacific lane factor)}
\]
\textit{[Note H1: 1.15 is calibrated for TPEB. For FEWB and TAWB, the factor is approximately
1.06 (shorter great-circle deviation on the Europe--Asia and transatlantic lanes). P1: replace
the scalar with a per-arc or per-trade-lane parameter $\gamma_{ij}$.]}

\paragraph{Transit time:}
\[
  \mu_{ij}^{\text{OC}} = \frac{\text{dist}_{ij}^{\text{sail}}}{600\,\text{km/day}}
  \quad (\approx 13.5\,\text{knots average, inclusive of port approach and anchorage})
\]

\paragraph{Validation:} SHA$\to$USLAX geodesic $\approx 9{,}200$\,km. Sailing distance
$\approx 10{,}580$\,km. Transit $\approx 17.6$\,days. Published carrier schedule for this
lane: 14--18 days. Estimate is within range.

\paragraph{Std dev:} $\sigma_{ij}^{\text{OC}} = 0.10\,\mu_{ij}^{\text{OC}}$ (10\% CV).

%%--------------------------------------------------------------------
\section{Instance Generator Specification}
\label{sec:instance}
%%--------------------------------------------------------------------

The instance generator produces synthetic but geographically and commercially realistic
problem instances for solver testing. The generator is demand-first: commodity requests
are generated first, then supply (sailings) is built to cover that demand at a configurable
capacity utilization level.

\textbf{Note on build process.} The generator implementation is a joint session between the
developer and the system designer. Do not implement independently.

\subsection{Configuration Schema}

All instance parameters are controlled by a \texttt{GeneratorConfig} object:

\begin{verbatim}
GeneratorConfig:

  # Trade lanes (drives POL/POD sets, carrier services, realistic schedules)
  trade_lanes: List[str]         # e.g. ["TPEB", "FEWB", "TAWB"]

  # Demand
  n_commodities: int             # total commodities to generate
  start_date: date               # scenario anchor date
  date_horizon_days: int         # window over which cargo-ready dates spread
  cargo_ready_spread_days: int   # subset of horizon for cargo-ready sampling

  # Commodity profile
  volume_range_cbm: [float, float]      # [min, max]
  weight_range_kg: [float, float]       # [min, max] — upper end tests weight-limited n_k
  deadline_slack_days: [int, int]       # [min, max] days added beyond fastest feasible path
  cargo_type_mix: dict                  # {general: 0.80, hazmat: 0.10, reefer: 0.10}
                                        # MVP: all routed as general; mix recorded only

  # Supply calibration
  load_factor: float             # target (total TEU demand) / (total TEU capacity)
                                 # 0.5 = easy, 0.8 = realistic tight, >1.0 = capacity-infeasible
  target_infeasible_fraction: float  # fraction of commodities with no feasible path
                                     # injected deliberately; tests infeasibility handling
  sailings_per_week_per_string: int  # 1 = weekly service (standard)
  n_carriers_per_lane: int           # competing carriers per trade lane

  # Carrier allocation (string-level)
  allocation_per_service_month_teu: int  # forwarder contracted block per string per month
  initial_utilization_fraction: float   # fraction already booked at solve time
                                        # rem = alloc * (1 - initial_utilization_fraction)
  target_alloc_utilization_fraction: float  # target fraction of alloc consumed at solve optimum
                                            # used to calibrate alloc vs cap (see step 5)
  vessel_cap_alpha: float               # alpha in cap_{ij} = alpha * alloc(s,t); default 0.20

  # Container cost (market-calibrated defaults)
  feu_base_rate_range_usd: [float, float]   # [min, max], e.g. [2000, 6000]
  teu_feu_cost_ratio: float                 # TEU rate = FEU rate * ratio; TPEB: ~0.79, FEWB: ~0.56, TAWB: ~0.78 (Xeneta 2020-2024)
  feu_hc_premium_usd: float                 # surcharge for 40'HC vs 40' standard; ~$50-150
  baf_fraction: float                       # BAF as fraction of base rate; ~0.15-0.25
  thc_pol_range_usd: [float, float]         # per container
  thc_pod_range_usd: [float, float]         # per container

  # Transit time
  transit_sigma_fraction: float   # sigma = fraction * mean for all arc types; default 0.12
                                  # [Note M5: Section 9 specifies 0.15 for trucking and 0.10
                                  #  for ocean. This config field is a single override. P1:
                                  #  split into trucking_sigma_fraction and ocean_sigma_fraction.]
  vessel_speed_kmday: float       # default 600 (~13.5 knots)

  # Random seed
  random_seed: int
\end{verbatim}

\subsection{Trade Lane Definitions}

\begin{center}
\begin{tabular}{lllll}
\toprule
Code & Direction & POL ports & POD ports & Example strings \\
\midrule
TPEB & Asia $\to$ US West & SHA, NGB, SZX & USLAX, USLGB, USSEA & MSC Tiger, CMA CGM AEX-1, ONE NE5 \\
FEWB & Asia $\to$ Europe & SHA, NGB, SZX, SGSIN & NLRTM, DEHAM, GBFXT, BEFBR & MSC Silk, Maersk AE-1, CMA CGM FAL1 \\
TAWB & Europe $\to$ US East & NLRTM, DEHAM, GBFXT, BEFBR & USNYK, USSAV, USBAL, USORF & MSC Britannia, Maersk TA1, Hapag Lloyd TAE \\
TPWB & US West $\to$ Asia & USLAX, USLGB, USSEA & SHA, NGB, SZX & (return legs of TPEB strings) \\
\bottomrule
\end{tabular}
\end{center}

\subsection{Demand-First Generation Algorithm}

\begin{enumerate}
  \item \textbf{Generate commodities.}
    \begin{enumerate}[noitemsep]
      \item Sample origin city from cities near the POLs of the configured trade lanes.
      \item Sample destination city from cities near the PODs.
      \item Sample $v_k \sim \text{Uniform}(v_{\min}, v_{\max})$,
            $w_k \sim \text{Uniform}(w_{\min}, w_{\max})$.
      \item Compute $n_k = \max(\ceil{v_k/76}, \ceil{w_k/26{,}500})$.
      \item Sample $t_k^{\text{rdy}}$ uniformly over \texttt{[start\_date, start\_date + cargo\_ready\_spread\_days]}.
    \end{enumerate}

  \item \textbf{Compute minimum feasible time for each commodity.}
    \begin{enumerate}[noitemsep]
      \item For each (origin, POL) pair: compute pre-carriage transit $\mu_{o(k),j}^{\text{PC}}$.
      \item Add CY buffer (4 days).
      \item Add fastest ocean transit to any reachable POD.
      \item Add port dwell.
      \item Add fastest inland transit to destination.
      \item $\text{min\_time}_k$ = minimum over all (origin $\to$ POL $\to$ POD $\to$ destination) paths.
    \end{enumerate}

  \item \textbf{Set deadlines.}
    \begin{itemize}[noitemsep]
      \item Normal commodities: $T_k^{\text{dead}} = t_k^{\text{rdy}} + \text{min\_time}_k
            + \text{Uniform}(\text{deadline\_slack\_min}, \text{deadline\_slack\_max})$.
      \item Infeasible commodities (fraction = \texttt{target\_infeasible\_fraction}):
            $T_k^{\text{dead}} = t_k^{\text{rdy}} + 0.7 \times \text{min\_time}_k$ (impossible).
    \end{itemize}

  \item \textbf{Generate sailings (demand-driven).}
    \begin{enumerate}[noitemsep]
      \item For each (POL, POD, trade lane) pair, determine the cargo-ready date window
            requiring coverage: $[t_k^{\text{rdy}}, t_k^{\text{rdy}} + \max\,\mu^{\text{PC}}]$ across all commodities.
      \item Generate weekly sailings (one per string per week) covering this window.
      \item Assign realistic carrier, service, and vessel names from the trade lane table.
      \item Compute $\mu_{ij}^{\text{OC}}$ and $\sigma_{ij}^{\text{OC}}$ via Section~\ref{sec:transit}.
      \item Set $\text{CYC}_{ij} = \text{ETD}_{ij} - 4\,\text{days}$.
            \textit{Sync note: this 4-day buffer must match the value used in step 2b of the
            minimum-feasible-time computation. Both occurrences must be updated together if
            per-port CYC variation is introduced (P1 extension).}
      \item \textbf{Calibrate per-sailing capacity to hit \texttt{load\_factor}:}
            $\text{cap}_{ij}^{\text{TEU}} = \text{round}(\text{total TEU demand on this lane}) /
            (\text{n\_sailings} \times \texttt{load\_factor})$.
            \textit{[Note C3: Circular dependency --- $\text{cap}_{ij}^{\text{TEU}}$ depends on
            total TEU demand, which depends on which sailings are feasible (step 4a/b), which
            depends on schedule coverage. Resolve by computing demand in step 1 first (with
            $n_k$ container counts), then setting cap. Do not re-run step 1 after setting cap.]}
    \end{enumerate}

  \item \textbf{Set allocation state.}
    $\text{rem}(s, t) = \texttt{allocation\_per\_service\_month\_teu}
    \times (1 - \texttt{initial\_utilization\_fraction})$.
    \textit{Calibration: set
    $\texttt{allocation\_per\_service\_month\_teu} \approx
    \Bigl(\sum_{(i,j):\, s_{ij}=s,\, t_{ij}=t} \text{cap}_{ij}^{\text{TEU}}\Bigr)
    \times \texttt{target\_alloc\_utilization\_fraction}$
    so that the string allocation constraint (P.3) binds at the intended load level.}

  \item \textbf{Pre-compute container mix.} For each (commodity, feasible sailing) pair, run
    the mix algorithm (Section~\ref{sec:mix-algo}) to get $f_k^*$, $t_k^*$, $\text{cost}(k,ij)$,
    $\text{slots}(k,ij)$.

  \item \textbf{Run subgraph construction} (Section~\ref{sec:prefilter}).
    Log infeasible commodities. Verify infeasible fraction is within 20\% of target;
    adjust \texttt{deadline\_slack} if needed and re-run.

  \item \textbf{Run decomposition} (Section~\ref{sec:decomp}). Log connected components and
    their sizes. This verifies that TPEB/FEWB mixed instances decompose as expected.

  \item \textbf{Serialize} to JSON with schema matching parameter tables above.
    Include random seed and config for reproducibility.
\end{enumerate}

\subsection{Solver Concerns}
\label{sec:solver}

\begin{itemize}[noitemsep]
  \item \textbf{Instance size:} At $|K|=50$, $\approx 72$ arcs, and $\approx 30$ feasible arcs per commodity, total binary variables $\approx 1{,}500$. Trivial for HiGHS.
  \item \textbf{Symmetry:} Inject small random perturbations ($\pm 1\%$) to FEU rates to break
        symmetry between near-identical sailings on the same lane.
  \item \textbf{LP bound:} The string allocation constraint (P.3) is the primary source of LP bound weakness. For tight-capacity instances, the gap may require B\&B. HiGHS handles this well for MVP sizes.
        \textit{[Note M4: The gap bound $\bigl(f_k^*\big|_{ij} + t_k^*\big|_{ij}\bigr) - n_k \leq 1$
        in §4.5 Remark is proven by a cost argument: $2\rho > 1$ for all $\rho > 0.5$, so the
        optimizer never replaces 1 FEU with 2 TEUs. The current remark text mentions ``three TEUs''
        which is the wrong comparison unit --- the binding case is 2 TEUs vs 1 FEU. Rephrase in
        a future edit.]}
\end{itemize}

%%--------------------------------------------------------------------
\section{Deferred: P1 Extensions}
\label{sec:deferred}
%%--------------------------------------------------------------------

\subsection{Commodity-Specific Customs Inspection Model (P1)}

Dwell arc weight becomes a function of commodity attributes:
\[
  \mu_k^{\text{DW}}(p) = \mu_p^{\text{unload}} + t^{\text{no-exam}}
    + p_k^{\text{exam}}(p) \cdot t^{\text{exam}}
\]
\[
  p_k^{\text{exam}}(p) = \operatorname{clip}\!\left(
    \lambda_p^{\text{base}} \cdot \eta_k^{\text{HS}} \cdot \eta_k^{\text{orig}} \cdot (1 - \delta_k^{\text{CTPAT}}),\; 0,\; 1
  \right)
\]

\subsection{Time-Phased Capacity Release (P1)}

Even with $\text{rem}(s,t)$ available, a forwarder may not want to commit all remaining
capacity in one routing run — future urgent shipments may need it. This requires
forecasting future demand and releasing capacity in tranches:
\[
  \text{open}(s,t,\ell) \leq \text{rem}(s,t)
\]
where $\ell$ is the routing run (batch) index and $\text{open}(s,t,\ell)$ is the capacity
released for run $\ell$. Requires a demand forecast model. Deferred.

\subsection{Probabilistic Delivery Constraint (P1)}

Replace deterministic deadline feasibility with:
\[
  \Pr\!\left(\text{delivery time}_k \leq T_k^{\text{dead}}\right) \geq \alpha_k
\]
where $\alpha_k \in \{0.80, 0.90, 0.95\}$ maps to service level (Economy/Standard/Express).

\subsection{Other P1 Items}

TEU-only container option for heavy cargo; multi-objective Pareto frontier (cost vs.\ transit time vs.\ on-time probability); port storage cost in objective; carrier allocation caps as soft constraints with penalty; transshipment arcs; multi-HS-code commodity schema (real FCL bookings can consolidate multiple products with different HS codes --- MVP uses a single $\text{hs}_k$ field, which is harmless while the customs model is deferred; P1 should replace $\text{hs}_k$ with a list and derive inspection risk as the maximum risk tier across all codes in the shipment).

%%--------------------------------------------------------------------
\section{Open Items}
%%--------------------------------------------------------------------

\begin{enumerate}
  \item \textbf{String allocation granularity:} Current model uses monthly periods for
        allocation. Some carrier contracts use 13 four-week periods per year rather than
        calendar months. Confirm with design partners before implementing.

  \item \textbf{Multi-POL / transshipment strings:} Some strings have calls at multiple
        Chinese ports (e.g., SHA $\to$ NGB $\to$ SZX before crossing). A booking on SHA$\to$USLAX
        and NGB$\to$USLAX on the same sailing are not independent — they're on the same vessel.
        The current model treats them as independent sailings, which may overstate available
        capacity. Deferred to when transshipment arcs are added in P1.

  \item \textbf{CYC compliance risk:} The subgraph construction uses mean transit times.
        A pre-carriage arc with mean transit exactly at the CYC has $\approx 50\%$ rollover risk.
        Flag this in the agent's output when the margin is $< 0.5$ days.

  \item \textbf{Infeasibility report schema:} Define the JSON schema for the infeasibility
        report before coding the pre-filter. Minimum fields: \texttt{commodity\_id, reason,
        earliest\_achievable\_days, deadline\_shortfall\_days, nearest\_feasible\_deadline}.

  \item \textbf{Generator: sailing schedule realism:} Real carrier schedules have port-specific
        cutoff windows, vessel-rotation identifiers, and multi-port calls. The generator
        should reproduce these from public schedule data, not fabricate them.

  \item \textbf{BSA contract unit convention:} The model uses TEU slots internally for all
        allocation quantities ($\text{alloc}(s,t)$, $\text{util}(s,t)$, $\text{rem}(s,t)$).
        Real BSA contracts on TPEB are commonly quoted in FEU; other trade lanes may use TEU
        or gross tonnage. Confirm with design partner: (a) are contracts quoted in FEU or TEU
        on the primary TPEB lanes, and (b) is the $\times 2$ conversion
        ($\text{alloc}_{\text{FEU}} \times 2 = \text{alloc in TEU slots}$) standard across
        all contract types? See also the unit convention remark in
        Section~\ref{sec:alloc-params}.

  \item \textbf{Vessel capacity proxy (P.2):} The per-sailing cap $\text{cap}_{ij}^{\text{TEU}}
        = \alpha \cdot \text{alloc}(s_{ij}, t_{ij})$ with $\alpha = 0.20$ is a placeholder.
        The right estimate is the fraction of the vessel's actual TEU capacity reserved for
        this forwarder on this string --- which requires carrier contract data or an AIS-based
        slot utilization feed. Candidate approaches: (a) request utilization reports directly
        from carriers in the BSA contract, (b) infer from AIS vessel identifiers and
        known vessel TEU capacities (publicly available), (c) use $\alpha = 0.5$ as a
        less conservative default once nominal carrier vessel sizes are known.
        \textit{Until real data is available, $\alpha = 0.20$ keeps the constraint non-binding
        while providing the structural guard against degenerate solutions.}
\end{enumerate}

\end{document}
